\pagestyle{empty}
\tight
\begin{tblr}{width=\textwidth,colspec={|Q[c,m,1.9cm]|X[l,m]|},hlines,vlines,hspan=minimal,row{1}={bg=headblue},rowsep=1pt,colsep=3pt}
\SetCell{c,m}\hfil\logo\hfil & \textbf{POLITEKNIK NEGERI MALANG}\par\textbf{JURUSAN TEKNOLOGI INFORMASI}\par\textbf{PROGRAM STUDI: D4 TEKNIK INFORMATIKA} \\
\SetCell[c=2]{c} \textbf{RENCANA PEMBELAJARAN SEMESTER (RPS)} & \\
\end{tblr}
\begin{longtblr}{width=\textwidth,colspec={|Q[l,m,1.9cm]|X[0.85,l,m]|X[1.45,l,m]|X[0.75,l,m]|X[1.15,l,m]|},hlines,vlines,hspan=minimal,rowsep=1pt,colsep=2pt,row{1,3}={bg=headgray,font=\bfseries}}
MATA KULIAH & KODE & BOBOT (SKS)/jam & SEMESTER & TGL. PENYUSUNAN \\
Jaringan Komputer & RTI254005 & 2 SKS/4 jam & 4 & 1 Juli 2025 \\
OTORISASI & Dosen Pengembang RPS & Koordinator RMK & \SetCell[c=2]{l} Ka PRODI &  \\
 & Vipkas Al Hadid Firdaus, S.T., M.T. &  & \SetCell[c=2]{l}{\vspace{8mm}\par Dr. Ely Setyo Astuti, S.T., M.T.} &  \\
\SetCell[r=12]{l} Capaian Pembelajaran (CP) & \SetCell[c=4]{l,bg=headgray,font=\bfseries} Capaian Pembelajaran Lulusan Program Studi (CPL-Prodi) &  &  &  \\
 & CPL05 & \SetCell[c=3]{l} Mampu menghasilkan solusi teknologi komputasi dan infrastruktur digital yang sesuai dengan kebutuhan organisasi dan pengguna. &  &  \\
 & CPL09 & \SetCell[c=3]{l} Mampu menjelaskan cara kerja sistem komputer dan menerapkan pendekatan komputasi untuk menyelesaikan masalah. &  &  \\
 & CPL10 & \SetCell[c=3]{l} Mampu menganalisis persoalan kompleks dengan wawasan multidisiplin untuk menghasilkan solusi di bidang informatika dan ilmu komputer. &  &  \\
 & \SetCell[c=4]{l,bg=headgray,font=\bfseries} Capaian Pembelajaran mata kuliah (CPL-MK) &  &  &  \\
 & CPMK0504 & \SetCell[c=3]{l} Mahasiswa mampu merancang arsitektur, memilih protokol, dan mengkonfigurasi infrastruktur jaringan komputer. &  &  \\
 & CPMK0901 & \SetCell[c=3]{l} Mahasiswa mampu menjelaskan cara kerja sistem komputer dan komponen utama infrastruktur jaringan. &  &  \\
 & CPMK1002 & \SetCell[c=3]{l} Mahasiswa mampu menganalisis mekanisme keamanan jaringan untuk melindungi data dan layanan digital. &  &  \\
 & \SetCell[c=4]{l,bg=headgray,font=\bfseries} Kemampuan akhir tiap tahapan belajar (Sub-CPMK) &  &  &  \\
 & SCPMK0504-03201 & \SetCell[c=3]{l} Mahasiswa mampu menganalisis arsitektur jaringan (OSI, TCP/IP) dan mengkonfigurasi protokol pada setiap lapisan. &  &  \\
 & SCPMK0901-03202 & \SetCell[c=3]{l} Mahasiswa mampu menjelaskan komponen jaringan (router, switch, media transmisi) dan cara kerja protokol komunikasi. &  &  \\
 & SCPMK1002-03203 & \SetCell[c=3]{l} Mahasiswa mampu menganalisis ancaman keamanan jaringan dan merekomendasikan mekanisme perlindungan yang sesuai. &  &  \\
\SetCell[r=2]{l} Penilaian & \SetCell[c=4]{l} *Nilai CPMK didapat dari agregasi nilai SUB-CPMK yang dibebankan &  &  &  \\
 & \SetCell[c=4]{l}{\tiny\begin{tblr}{width=\linewidth,colspec={|Q[c,m,0.1\linewidth]|Q[c,m,0.16\linewidth]|X[l,m]|X[l,m]|X[l,m]|X[l,m]|X[l,m]|X[l,m]|Q[c,m,0.08\linewidth]|},hlines,vlines,hspan=minimal,rowsep=1pt,colsep=0.7pt,row{1,2}={bg=headgray,font=\bfseries}}
Kode CPMK & Kode SCPMK & \SetCell[c=6]{c} Bobot per Bentuk Penilaian & & & & & & Total Bobot \\
 & & Tugas & Kuis 1 & UTS & Kuis 2 & PBL & UAS & \\
CPMK0504 & SCPMK0504-03201 & 12\%: analisis arsitektur jaringan & 0\% & 5\%: protokol jaringan & 0\% & 12\%: desain jaringan & 12\%: presentasi desain jaringan & 41\% \\
CPMK0901 & SCPMK0901-03202 & 10\%: komponen dan protokol jaringan & 5\%: model OSI dan TCP/IP & 8\%: komponen jaringan & 0\% & 0\% & 0\% & 23\% \\
CPMK1002 & SCPMK1002-03203 & 8\%: analisis keamanan jaringan & 0\% & 0\% & 5\%: keamanan jaringan dasar & 10\%: analisis keamanan & 13\%: analisis keamanan jaringan & 36\% \\
\SetCell[c=8]{r}\textbf{Total Per Penilaian} & & & & & & & & \textbf{100\%} \\
\end{tblr}} &  &  &  \\
Diskripsi Singkat Mata Kuliah & \SetCell[c=4]{l} Jaringan Komputer membangun pemahaman komprehensif tentang arsitektur jaringan, protokol komunikasi, perangkat jaringan, dan keamanan jaringan untuk merancang infrastruktur jaringan yang efektif dan aman. &  &  &  \\
Bahan Kajian / Materi Pembelajaran & \SetCell[c=4]{l}\begin{enumerate}\item Pengantar jaringan, model OSI, TCP/IP, dan media transmisi\item Protokol application, transport, dan network layer\item Wireless networking dan data link layer\item Keamanan jaringan: ancaman, firewall, VPN, dan monitoring\end{enumerate} &  &  &  \\
\SetCell[r=2]{l} Pustaka & \SetCell[c=4]{l} \textbf{Utama:}\begin{enumerate}\item Tanenbaum, A. 2021. Computer Networks (6th ed.). Pearson.\item Forouzan, B. 2018. Data Communications and Networking. McGraw-Hill.\item Kurose, J. \& Ross, K. 2021. Computer Networking: A Top-Down Approach. Pearson.\end{enumerate} &  &  &  \\
 & \SetCell[c=4]{l} \textbf{Pendukung:} Materi LMS, dokumentasi resmi, dan jobsheet praktikum. &  &  &  \\
Media Pembelajaran & \SetCell[c=2]{l} \textbf{Software:} Cisco Packet Tracer, Wireshark, LMS. &  & \SetCell[c=2]{l} \textbf{Hardware:} Komputer/laptop, proyektor. &  \\
Nama Dosen Pengampu & \SetCell[c=4]{l} Tim Pengajar Program Studi D4 TI &  &  &  \\
Matakuliah Syarat & \SetCell[c=4]{l} - &  &  &  \\
\end{longtblr}
\begin{longtblr}{width=\textwidth,colspec={|Q[c,m,0.06\textwidth]|X[1.35,l,m]|X[1.08,l,m]|X[1.16,l,m]|X[0.7,l,m]|X[1.24,l,m]|X[1.06,l,m]|X[0.98,l,m]|Q[c,m,0.05\textwidth]|},hlines,vlines,hspan=minimal,rowhead=2,row{1,2}={bg=headgreen,font=\bfseries},rowsep=1pt,colsep=1.5pt}
Mgg.\par Ke & Kemampuan Akhir Yang Direncanakan (Sub-CP-MK) & Bahan kajian (Materi Pembelajaran) & Bentuk dan Metode Pembelajaran & Estimasi Waktu & Pengalaman Belajar Mahasiswa & Kriteria \& Bentuk Penilaian & Indikator Penilaian & Bobot Penilaian (\%) \\
(1) & (2) & (3) & (4) & (5) & (6) & (7) & (8) & (9) \\
1 & SCPMK0901-03202 & Pengantar jaringan komputer: definisi, tipe, dan topologi. & Modalitas: Blended Learning. Bentuk: Luring/Daring. Metode: Case Method, praktik konfigurasi jaringan, studi kasus, dan peer review. & 1 x 2 x 50' tatap muka; 1 x 2 x 50' tugas/praktik mandiri. & Mahasiswa mengerjakan aktivitas terarah untuk pengantar jaringan komputer: definisi, tipe, dan topologi dan menunjukkan evidence hasil belajar. & Bentuk penilaian: Pengantar Jaringan Komputer: Definisi, Tipe, dan Topologi. Mengacu rubrik RTM. & Ketepatan konsep; kualitas konfigurasi/analisis; validasi dan dokumentasi. & 2 \\
2 & SCPMK0901-03202 & Model OSI: 7 lapisan dan fungsinya. & Modalitas: Blended Learning. Bentuk: Luring/Daring. Metode: Case Method, praktik konfigurasi jaringan, studi kasus, dan peer review. & 1 x 2 x 50' tatap muka; 1 x 2 x 50' tugas/praktik mandiri. & Mahasiswa mengerjakan aktivitas terarah untuk model OSI: 7 lapisan dan fungsinya dan menunjukkan evidence hasil belajar. & Bentuk penilaian: Model OSI: 7 Lapisan dan Fungsinya. Mengacu rubrik RTM. & Ketepatan konsep; kualitas konfigurasi/analisis; validasi dan dokumentasi. & 2 \\
3 & SCPMK0901-03202 & TCP/IP suite dan protokol internet. & Modalitas: Blended Learning. Bentuk: Luring/Daring. Metode: Case Method, praktik konfigurasi jaringan, studi kasus, dan peer review. & 1 x 2 x 50' tatap muka; 1 x 2 x 50' tugas/praktik mandiri. & Mahasiswa mengerjakan aktivitas terarah untuk TCP/IP suite dan protokol internet dan menunjukkan evidence hasil belajar. & Bentuk penilaian: TCP/IP Suite dan Protokol Internet. Mengacu rubrik RTM. & Ketepatan konsep; kualitas konfigurasi/analisis; validasi dan dokumentasi. & 3 \\
4 & SCPMK0901-03202 & Media transmisi dan perangkat jaringan. & Modalitas: Blended Learning. Bentuk: Luring/Daring. Metode: Case Method, praktik konfigurasi jaringan, studi kasus, dan peer review. & 1 x 2 x 50' tatap muka; 1 x 2 x 50' tugas/praktik mandiri. & Mahasiswa mengerjakan aktivitas terarah untuk media transmisi dan perangkat jaringan dan menunjukkan evidence hasil belajar. & Bentuk penilaian: Media Transmisi dan Perangkat Jaringan. Mengacu rubrik RTM. & Ketepatan konsep; kualitas konfigurasi/analisis; validasi dan dokumentasi. & 2 \\
5 & SCPMK0901-03202 & Kuis 1: model OSI dan TCP/IP. & Modalitas: Blended Learning. Bentuk: Luring/Daring. Metode: Case Method, praktik konfigurasi jaringan, studi kasus, dan peer review. & 1 x 2 x 50' tatap muka; 1 x 2 x 50' tugas/praktik mandiri. & Mahasiswa mengerjakan aktivitas terarah untuk kuis 1: model OSI dan TCP/IP dan menunjukkan evidence hasil belajar. & Bentuk penilaian: Kuis 1: Model OSI dan TCP/IP. Mengacu rubrik RTM. & Ketepatan konsep; kualitas konfigurasi/analisis; validasi dan dokumentasi. & 5 \\
6 & SCPMK0504-03201 & Application layer: HTTP, FTP, DNS, SMTP. & Modalitas: Blended Learning. Bentuk: Luring/Daring. Metode: Case Method, praktik konfigurasi jaringan, studi kasus, dan peer review. & 1 x 2 x 50' tatap muka; 1 x 2 x 50' tugas/praktik mandiri. & Mahasiswa mengerjakan aktivitas terarah untuk application layer: HTTP, FTP, DNS, SMTP dan menunjukkan evidence hasil belajar. & Bentuk penilaian: Application Layer: HTTP, FTP, DNS, SMTP. Mengacu rubrik RTM. & Ketepatan konsep; kualitas konfigurasi/analisis; validasi dan dokumentasi. & 3 \\
7 & SCPMK0504-03201 & Transport layer: TCP vs UDP. & Modalitas: Blended Learning. Bentuk: Luring/Daring. Metode: Case Method, praktik konfigurasi jaringan, studi kasus, dan peer review. & 1 x 2 x 50' tatap muka; 1 x 2 x 50' tugas/praktik mandiri. & Mahasiswa mengerjakan aktivitas terarah untuk transport layer: TCP vs UDP dan menunjukkan evidence hasil belajar. & Bentuk penilaian: Transport Layer: TCP vs UDP. Mengacu rubrik RTM. & Ketepatan konsep; kualitas konfigurasi/analisis; validasi dan dokumentasi. & 3 \\
8 & SCPMK0901-03202, SCPMK0504-03201 & UTS: protokol dan lapisan jaringan. & Modalitas: Blended Learning. Bentuk: Luring/Daring. Metode: Case Method, praktik konfigurasi jaringan, studi kasus, dan peer review. & 1 x 2 x 50' tatap muka; 1 x 2 x 50' tugas/praktik mandiri. & Mahasiswa mengerjakan aktivitas terarah untuk UTS: protokol dan lapisan jaringan dan menunjukkan evidence hasil belajar. & Bentuk penilaian: UTS: Protokol dan Lapisan Jaringan. Mengacu rubrik RTM. & Ketepatan konsep; kualitas konfigurasi/analisis; validasi dan dokumentasi. & 13 \\
9 & SCPMK0504-03201 & Network layer: IP addressing dan routing. & Modalitas: Blended Learning. Bentuk: Luring/Daring. Metode: Case Method, praktik konfigurasi jaringan, studi kasus, dan peer review. & 1 x 2 x 50' tatap muka; 1 x 2 x 50' tugas/praktik mandiri. & Mahasiswa mengerjakan aktivitas terarah untuk network layer: IP addressing dan routing dan menunjukkan evidence hasil belajar. & Bentuk penilaian: Network Layer: IP Addressing dan Routing. Mengacu rubrik RTM. & Ketepatan konsep; kualitas konfigurasi/analisis; validasi dan dokumentasi. & 3 \\
10 & SCPMK0504-03201 & Subnetting dan CIDR. & Modalitas: Blended Learning. Bentuk: Luring/Daring. Metode: Case Method, praktik konfigurasi jaringan, studi kasus, dan peer review. & 1 x 2 x 50' tatap muka; 1 x 2 x 50' tugas/praktik mandiri. & Mahasiswa mengerjakan aktivitas terarah untuk subnetting dan CIDR dan menunjukkan evidence hasil belajar. & Bentuk penilaian: Subnetting dan CIDR. Mengacu rubrik RTM. & Ketepatan konsep; kualitas konfigurasi/analisis; validasi dan dokumentasi. & 3 \\
11 & SCPMK0504-03201 & Data link dan physical layer. & Modalitas: Blended Learning. Bentuk: Luring/Daring. Metode: Case Method, praktik konfigurasi jaringan, studi kasus, dan peer review. & 1 x 2 x 50' tatap muka; 1 x 2 x 50' tugas/praktik mandiri. & Mahasiswa mengerjakan aktivitas terarah untuk data link dan physical layer dan menunjukkan evidence hasil belajar. & Bentuk penilaian: Data Link dan Physical Layer. Mengacu rubrik RTM. & Ketepatan konsep; kualitas konfigurasi/analisis; validasi dan dokumentasi. & 3 \\
12 & SCPMK0504-03201 & Wireless networking: 802.11 dan Wi-Fi standards. & Modalitas: Blended Learning. Bentuk: Luring/Daring. Metode: Case Method, praktik konfigurasi jaringan, studi kasus, dan peer review. & 1 x 2 x 50' tatap muka; 1 x 2 x 50' tugas/praktik mandiri. & Mahasiswa mengerjakan aktivitas terarah untuk wireless networking: 802.11 dan Wi-Fi standards dan menunjukkan evidence hasil belajar. & Bentuk penilaian: Wireless Networking: 802.11 dan Wi-Fi Standards. Mengacu rubrik RTM. & Ketepatan konsep; kualitas konfigurasi/analisis; validasi dan dokumentasi. & 3 \\
13 & SCPMK1002-03203 & Kuis 2: keamanan jaringan dasar. & Modalitas: Blended Learning. Bentuk: Luring/Daring. Metode: Case Method, praktik konfigurasi jaringan, studi kasus, dan peer review. & 1 x 2 x 50' tatap muka; 1 x 2 x 50' tugas/praktik mandiri. & Mahasiswa mengerjakan aktivitas terarah untuk kuis 2: keamanan jaringan dasar dan menunjukkan evidence hasil belajar. & Bentuk penilaian: Kuis 2: Keamanan Jaringan Dasar. Mengacu rubrik RTM. & Ketepatan konsep; kualitas konfigurasi/analisis; validasi dan dokumentasi. & 5 \\
14 & SCPMK1002-03203 & Network security: ancaman, serangan, dan mitigasi. & Modalitas: Blended Learning. Bentuk: Luring/Daring. Metode: Case Method, praktik konfigurasi jaringan, studi kasus, dan peer review. & 1 x 2 x 50' tatap muka; 1 x 2 x 50' tugas/praktik mandiri. & Mahasiswa mengerjakan aktivitas terarah untuk network security: ancaman, serangan, dan mitigasi dan menunjukkan evidence hasil belajar. & Bentuk penilaian: Network Security: Ancaman, Serangan, dan Mitigasi. Mengacu rubrik RTM. & Ketepatan konsep; kualitas konfigurasi/analisis; validasi dan dokumentasi. & 3 \\
15 & SCPMK1002-03203 & Firewall, VPN, dan network monitoring. & Modalitas: Blended Learning. Bentuk: Luring/Daring. Metode: Case Method, praktik konfigurasi jaringan, studi kasus, dan peer review. & 1 x 2 x 50' tatap muka; 1 x 2 x 50' tugas/praktik mandiri. & Mahasiswa mengerjakan aktivitas terarah untuk firewall, VPN, dan network monitoring dan menunjukkan evidence hasil belajar. & Bentuk penilaian: Firewall, VPN, dan Network Monitoring. Mengacu rubrik RTM. & Ketepatan konsep; kualitas konfigurasi/analisis; validasi dan dokumentasi. & 3 \\
16 & SCPMK0504-03201, SCPMK1002-03203 & PBL: desain dan analisis jaringan. & Modalitas: Blended Learning. Bentuk: Luring/Daring. Metode: Case Method, praktik konfigurasi jaringan, studi kasus, dan peer review. & 1 x 2 x 50' tatap muka; 1 x 2 x 50' tugas/praktik mandiri. & Mahasiswa mengerjakan aktivitas terarah untuk PBL: desain dan analisis jaringan dan menunjukkan evidence hasil belajar. & Bentuk penilaian: PBL: Desain dan Analisis Jaringan. Mengacu rubrik RTM. & Ketepatan konsep; kualitas konfigurasi/analisis; validasi dan dokumentasi. & 22 \\
17 & SCPMK0504-03201, SCPMK1002-03203 & UAS: presentasi desain jaringan dan analisis keamanan. & Modalitas: Blended Learning. Bentuk: Luring/Daring. Metode: Case Method, praktik konfigurasi jaringan, studi kasus, dan peer review. & 1 x 2 x 50' tatap muka; 1 x 2 x 50' tugas/praktik mandiri. & Mahasiswa mengerjakan aktivitas terarah untuk UAS: presentasi desain jaringan dan analisis keamanan dan menunjukkan evidence hasil belajar. & Bentuk penilaian: UAS: Presentasi Desain Jaringan dan Analisis Keamanan. Mengacu rubrik RTM. & Ketepatan konsep; kualitas konfigurasi/analisis; validasi dan dokumentasi. & 25 \\
\end{longtblr}
